%%=============================================================================
%% CAPITULO 3 -- PROCEDIMENTOS METODOLOGICOS
%% Texto dummy do exercicio ia-6.1.
%%=============================================================================

\section{Procedimentos metodológicos}

A pesquisa é aplicada e de abordagem quantitativa. Um conjunto de perfis
sintéticos de investidores, variando objetivo, prazo e tolerância a risco, será
submetido a um protótipo de recomendador que escolhe entre os produtos da
Tabela~\ref{tab:produtos} e gera, para cada escolha, uma justificativa em
português.

Cada par formado por recomendação e justificativa será avaliado em duas
dimensões: a adequação do produto ao perfil, por regra explícita derivada da
regulação de suitability, e a fidedignidade do texto, pela comparação de cada
afirmação com a ficha do produto, classificando as divergências segundo a
tipologia de alucinação adotada no Capítulo~2.

\begin{table}[htb]
  \caption{Produtos de renda fixa cobertos pelo protótipo}
  \label{tab:produtos}
  \centering
  \begin{tabular}{lll}
    \toprule
    \textbf{Produto} & \textbf{Garantia} & \textbf{Imposto de renda (PF)} \\
    \midrule
    Tesouro Direto & Tesouro Nacional & tabela regressiva \\
    CDB            & FGC              & tabela regressiva \\
    LCI e LCA      & FGC              & isento            \\
    \bottomrule
  \end{tabular}
  \fonte{Elaborado pela autora (2026).}
\end{table}
